\documentclass[UTF8,11pt,a4paper]{ctexart}
\usepackage[margin=1.9cm,headheight=15pt]{geometry}
\usepackage{amsmath,amssymb,booktabs,array,graphicx,tikz,xcolor,hyperref,xurl,fancyhdr,enumitem}
\usetikzlibrary{arrows.meta,positioning}
\setmainfont{Arial}\setsansfont{Arial}\setmonofont{Menlo}
\setCJKmainfont{Songti SC}\setCJKsansfont{Heiti SC}
\definecolor{ink}{HTML}{163C50}\definecolor{accent}{HTML}{007C91}
\hypersetup{colorlinks=true,urlcolor=accent,linkcolor=ink,pdftitle={S85与S86：从历史投影到生成强对照},pdfauthor={AI辅助整理；依据项目实际记录}}
\renewcommand{\UrlFont}{\fontspec{Arial Unicode MS}\small}
\setlength{\parskip}{.48em}\setlength{\parindent}{1.5em}\linespread{1.06}
\setlength{\emergencystretch}{2em}\renewcommand{\arraystretch}{1.13}
\pagestyle{fancy}\fancyhf{}\lhead{\small S85 / S86：零基础教学增补}\rhead{\small 科学截点 09-11 05:34}\cfoot{\small\thepage}
\ctexset{section={format=\Large\bfseries\sffamily\color{ink},beforeskip=0pt,afterskip=.7em}}
\newcommand{\teaching}[1]{\begin{center}\fcolorbox{accent}{accent!5}{\parbox{.91\linewidth}{\small #1}}\end{center}}
\begin{document}


\section*{1. proposal主线，本轮补的是哪一段？}

\begin{center}{\large\bfseries 从历史照片投影，到真正影响生成}\end{center}

\teaching{\textbf{科学状态截点：北京时间2026-09-11 05:34:29。}S85已完成并通过不同作者复算；S86已正式启动，监督回执仅确认G0完成9/50步、仍在运行。四臂尚未完成、未评分。本增补不覆盖旧142页汇报，也不预写成功。}

研究最终希望回答：模型走过一个地方、看过若干画面以后，能否在未来视角仍生成结构合理、与过去相容的场景？这里的“记住”不是单纯把照片存到硬盘，而是后续生成真的使用了适当的历史信息，而且这种使用带来可独立评价的收益。
\par


本轮把路线拆成五个连续动作：存下历史；选择要用的历史；把历史内容定位到请求视角；让生成器消费它；用真实参考检查是否有益。前一动作成功不会自动保证后一动作成功。例如，历史照片真实、投影程序正确，也可能因为深度预测错误而提供错误位置；强行使用它还可能让生成更差。
\par


S82曾用四张历史照片运行已有CUT3R模型，保存预测几何；S83固定给定相机和内参，做100步拟合，保存清理前状态；S85把这份状态和历史颜色投到四个目标相机，已经形成可复查的warp；S86才把同一warp接到原VMem生成过程，比较四种使用方式。这里的warp就是“历史照片按预测几何搬到新视角后的参考图”。
\par


\begin{center}\begin{tikzpicture}[>=Latex,font=\small,box/.style={draw=accent,fill=accent!5,rounded corners=2pt,align=center,text width=2.55cm,minimum height=1.2cm}]
            \node[box](a){历史照片\\S82预测几何};\node[box,right=4mm of a](b){固定相机\\S83拟合深度};
            \node[box,right=4mm of b](c){搬到目标视角\\S85已核投影};\node[box,right=4mm of c](d){四种生成对照\\S86运行中};
            \draw[->](a)--(b);\draw[->](b)--(c);\draw[->](c)--(d);\end{tikzpicture}\end{center}
            \noindent\small\textcolor{accent}{流程教学图；末框按本报告截点仍未完成。}\normalsize

这些是普通基线与机制区分的实际推进。还缺什么同样明确：一组相关视角不能验证长程、动态物体或跨场景泛化；普通投影、融合和终端对照也不是新方法本身。当前仍为new\_method\_validated=false、novelty\_authorization=NONE。
\par


阅读顺序：先看第2、3页的手算；第4、5页看真实投影与完整分母；第6、7页理解四臂对照；第8、9页理解评分与否决；第10页准备汇报。
\par


\clearpage

\section*{2. 一个像素怎样成为另一个视角的像素？}

一张RGB照片中，像素(u,v)只直接存红、绿、蓝三个颜色值。相机内参K提供焦距fx、fy和主点cx、cy，用来把像素位置变成一条空间射线。深度Z决定这条射线上的具体位置。此处Z是沿相机光轴的深度，不是相机到物体的直线距离range。
\par


针孔相机公式为：X=(u-cx)Z/fx，Y=(v-cy)Z/fy，第三个坐标就是Z。相机外参P=(R,c)说明相机在世界中的朝向R和位置c，所以世界点为R(X,Y,Z)+c。再减去目标相机位置、按目标朝向转换回来，最后做u′=fx′X′/Z′+cx′、v′=fy′Y′/Z′+cy′。
\par


\teaching{\textbf{针孔公式教学表达：}
    \[\mathbf{X}_s=Z\begin{bmatrix}(u-c_x)/f_x\\(v-c_y)/f_y\\1\end{bmatrix},\quad
    \mathbf{X}_t=R_t^T(R_s\mathbf{X}_s+c_s-c_t),\quad
    u'=f'_x (\mathbf{X}_t)_x/(\mathbf{X}_t)_z+c'_x.\]
    Z单位是米，焦距单位是像素；坐标方向与每个变换的输入输出必须对应。}

\noindent\textbf{【以下全部是人工教学数值，不是本次实测。】}设fx=fy=100像素、主点(50,50)，像素(60,50)的Z=2米。先有X=(60-50)×2/100=0.2米，Y=0，得到(0.2,0,2)。它到相机中心的距离是sqrt(0.2²+2²)，约2.010米，说明range与Z并不相同。
\par


再设目标相机向右移动0.1米、没有旋转，目标坐标变成(0.1,0,2)。投回图像后，横坐标是55（100×0.1÷2+50），纵坐标仍为50：相机向右走，固定物体在图像里向左移动5像素。这是坐标变换带来的结果，不是模型“想象”了一次运动。
\par


如果Z估错、K不准或相机位置有误，最后像素位置也可能错；两份程序都实现同一公式，不意味着输入几何真实。S85沿用近似内参和光学相机约定，没有偷偷加尺度拟合或把坐标轴翻转来改善结果。
\par


\clearpage

\section*{3. 为什么一个点有四个候选，还会出现孔洞？}

投影结果通常不是整数像素位置。以人工点(10.2,20.3)为例，它附近的四个整数像素是(10,20)、(11,20)、(10,21)、(11,21)。横向小数为0.2，纵向小数为0.3，四个双线性权重依次为0.8×0.7=0.56、0.2×0.7=0.14、0.8×0.3=0.24、0.2×0.3=0.06。这一小例只是教学，不是抽取真实图上的测量。
\par


S85用“权重大于零”决定某个整数像素是否得到该点的候选，不用这些权重平均颜色。如果两个来源都到同一目标像素，先选目标相机Z较小者；Z完全相同才按历史19、18、13、12，再按源像素编号决胜。最后直接复制赢家的原RGB。第二候选是完整排序的第二项，可以与赢家Z相等，不能称为第二层真实表面。
\par


\begin{center}\small
    \begin{tabular}{lrrrr}\toprule
    人工目标格点&(10,20)&(11,20)&(10,21)&(11,21)\\\midrule
    正权重&0.56&0.14&0.24&0.06\\
    作用&候选资格&候选资格&候选资格&候选资格\\\bottomrule
    \end{tabular}\end{center}

再用人工例说明：红色候选在Z=2米，蓝色候选在Z=3米，本实现选红色。它没有证明红色真的可见，只说明在这份预测深度下红色排在前面。如果Z输入本来错了，程序也可能十分确定地选错。
\par


实际规则保留六类互斥状态：有效；源深度非有限；源深度非正；目标坐标非有限；目标在相机后方或Z为零；连续投影中心出画。目标Z非正时不能先做透视除法。576×576目标的连续中心必须在闭区间[0,575]²，不把出画中心强行夹回边缘，不用epsilon把小权重改成零。
\par


没有候选的目标像素叫孔洞。存储RGB可以为零，但独立mask才说明是否缺历史支持。真实黑色桌腿也可能RGB接近零，因此“非黑就是有效”的判断不成立。
\par


\clearpage

\section*{4. 四张真实保存的投影图，应当怎样看？}

以下是S85已导出并验收的四个目标投影。颜色来自历史实拍照片，位置由预测几何决定。灰格是标注的孔洞，不是原场景纹理。这些图不是目标视角实拍，也不是S86新生成。
\par


\begin{center}
\begin{minipage}{.46\linewidth}\centering\includegraphics[width=6.2cm]{images/target_20_holes_marked.png}\\{\small 目标 20：历史颜色投影}\end{minipage}\hfill
\begin{minipage}{.46\linewidth}\centering\includegraphics[width=6.2cm]{images/target_21_holes_marked.png}\\{\small 目标 21：历史颜色投影}\end{minipage}\\[3mm]
\begin{minipage}{.46\linewidth}\centering\includegraphics[width=6.2cm]{images/target_22_holes_marked.png}\\{\small 目标 22：历史颜色投影}\end{minipage}\hfill
\begin{minipage}{.46\linewidth}\centering\includegraphics[width=6.2cm]{images/target_23_holes_marked.png}\\{\small 目标 23：历史颜色投影}\end{minipage}\\[3mm]
\end{center}

\begin{center}\small\begin{tabular}{lrrrr}\toprule
目标 & 固定目标像素 & 有投影支持 & 孔洞 & 预测覆盖\\
\midrule
20 & 331776 & 312396 & 19380 & 94.1587\%\\
21 & 331776 & 292217 & 39559 & 88.0766\%\\
22 & 331776 & 267572 & 64204 & 80.6484\%\\
23 & 331776 & 263594 & 68182 & 79.4494\%\\
\bottomrule\end{tabular}\end{center}

目标20覆盖约94\%，只表示约94\%的目标格点存在至少一个预测候选，不能解释成准确率94\%。看见连续结构也不能直接断言深度正确；看见缺口也不能单凭彩图确定是遮挡、深度误差还是视角范围造成。
\par


S85投影真实运行2.293416秒，监督全过程2.378199秒；后续不同作者从原输入复算1.823490秒。短耗时来自已存数组上的投影和核验，不包含重新训练、生成视频或本次S86的两条完整链。
\par


\clearpage

\section*{5. 16对全部保留，数字到底在数什么？}

四张历史图各384×512=196608个像素；四图合计786432个独特的“历史帧ID＋像素ID”位置。每张再分别投到四个目标，所以有4×4=16对、3145728条源点记录。同一个历史像素会在四个目标里重复出现，这些行不是独立实验样本。
\par


下表保留全部16对，每行源点分母均为196608；数值直接整理自已接受的SUMMARY。本批源/目标非有限与非正四类状态均为0，出画源点保留。足迹数与源点数不是同一分母。
\par


\begin{center}\small\begin{tabular}{rrrrrr}\toprule
目标&历史&有效源点&出画源点&合格足迹&赢家足迹\\\midrule
20 & 12 & 139,539 & 57,069 & 558,156 & 140,293\\
20 & 13 & 153,689 & 42,919 & 614,756 & 67,726\\
20 & 18 & 137,344 & 59,264 & 549,376 & 50,000\\
20 & 19 & 137,754 & 58,854 & 551,016 & 54,377\\
\addlinespace[.3em]
21 & 12 & 141,992 & 54,616 & 567,968 & 149,050\\
21 & 13 & 141,666 & 54,942 & 566,664 & 58,708\\
21 & 18 & 124,932 & 71,676 & 499,728 & 44,703\\
21 & 19 & 125,324 & 71,284 & 501,296 & 39,756\\
\addlinespace[.3em]
22 & 12 & 123,450 & 73,158 & 493,800 & 147,587\\
22 & 13 & 113,666 & 82,942 & 454,664 & 49,314\\
22 & 18 & 100,145 & 96,463 & 400,580 & 35,144\\
22 & 19 & 102,866 & 93,742 & 411,464 & 35,527\\
\addlinespace[.3em]
23 & 12 & 108,438 & 88,170 & 433,752 & 150,815\\
23 & 13 & 95,867 & 100,741 & 383,468 & 43,776\\
23 & 18 & 85,961 & 110,647 & 343,844 & 33,233\\
23 & 19 & 89,594 & 107,014 & 358,376 & 35,770\\
\bottomrule\end{tabular}\end{center}

每个目标的赢家足迹总和等于该目标有支持的像素数，每个目标像素最多有一个赢家。目标20至23至少两候选像素分别为299445、271712、237796、225874；这些竞争也可来自同一张历史图的相邻像素，不必是跨帧矛盾或真实遮挡。
\par


独立核验读取2份原输入、1份消费副本和4份目标档案，核完整候选无重复排列、排序和全部摘要。120次compare是8个输入副本字段＋4×28次目标比较；四目标各有31字段结构检查，3个候选身份字段通过排列与排序联合验证。28个浮点数值比较项最大差均为0。它证明实现与保存一致，不是与世界真实几何的误差为零。
\par


\clearpage

\section*{6. 四种对照，各自回答一个问题}

G0：保持原VMem生成链，生成过程不融合warp。它回答“原方法在共同条件下本来怎样”。
\par


Gpaste：复用G0生成结果，在最终RGB画面中按图像mask与warp做λ=0.25的合成。它回答“如果只是最后贴回一些历史颜色，能解释多少变化”。没有额外denoiser链。
\par


Gterminal：复用G0真实最后一步保存量，在模型内部压缩表示latent里融合同一个warp，只重演最后Euler算术，然后解码。它回答“最后一次内部融合和解码是否已经足够”。它与RGB贴图不同，但也不需要再跑一整条denoiser链。
\par


Gguide：另一条完整原生成链，在每一步原CFG合并之后、Euler更新之前，都用同一λ=0.25融合。总计50步，后续状态自然重新计算。它回答“较早介入及其后续传播是否带来额外作用”。
\par


只比较Gguide与Gpaste不够：二者差别既有介入时点，也有RGB/latent域和解码方式。加入Gterminal后，Gguide与它共享末步warp、mask和强度，较早融合的作用才有更合适的强对照。即使胜出，50次累计干预剂量仍可能解释差别，不能自动叫新的几何推理能力。
\par


四臂共用四历史、四目标、同一warp与固定条件。两条完整链恢复相同的实际随机状态，并逐步核随机流；不是只给相同seed就算公平。模型内的后续预测允许改变，不能把干预后代强行冻结回G0。
\par


本次保留原VMem采样流程与权重，但VAE使用已声明的stabilityai/sd-vae-ft-mse变体；原SD2.1 VAE身份仍未确认，因此不声称完整原始系统的精确复现。固定运行设置为CPU FP32、8线程、50步、576×576，四历史与四目标共同组成8帧槽位。
\par


\clearpage

\section*{7. λ=0.25究竟改了多少？}

\noindent\textbf{【本页数值全部是人工教学，不是S86结果。】}令当前clean预测某一元素d=0.2，历史warp编码对应值g=0.8，覆盖比例m=1，历史保护标记h=0。实际权重w=λm(1-h)=0.25，融合后d\_g=(1-w)d+wg=0.75×0.2+0.25×0.8=0.35。变化是0.15，等于w(g-d)。
\par


若同一个8×8块只有一半图像位置有预测支持，avg\_pool得到m=0.5，则w=0.125，结果变为0.275。m是64个格点中有支持的比例，不是“这个位置正确的概率”。若为历史槽h=1或无支持m=0，则w=0，原值保持不变。
\par


\teaching{\textbf{人工教学，非实验数值：}
    \[w=\lambda m(1-h),\qquad d_g=(1-w)d+wg.\]
    \[0.75\times0.2+0.25\times0.8=0.35,\qquad
      0.875\times0.2+0.125\times0.8=0.275.\]
    第一式覆盖比例为1；第二式为0.5。历史保护或无支持时，直接保留原值。}

S86实际将S85黑洞RGB0变成编码输入-1，按已有VAE的确定性均值编码并乘0.18215；不是随机抽一次latent。VAE感受野跨越多个像素，所以填洞方式可能影响邻近latent，不能把72×72的软mask当严格局部物理可靠性。
\par


原Euler同一步可在精确算术下写成E=a·x\_tilde+(1-a)d，a=sigma\_next/sigma\_hat。若教学设x\_tilde=1、a=0.5，则原输出0.6，融合后输出0.675，变化0.075=(1-a)×0.25×(0.8-0.2)。
\par


最后a=0时，输出在精确算术下等于本步的clean预测。但真实FP32有运算次序和相消误差，所以Gterminal仍按原to\_d、dt、x+dt×derivative重演，不能为了简便直接把融合clean当逐字节相同的最终latent。多步Gguide此前已改变当前预测d，末步等式也不意味着它与从G0派生的Gterminal完全相同。
\par


\clearpage

\section*{8. MSE的完整手算与本次真实分母}

MSE是平均平方误差：逐项求“预测-参考”的差、平方、加总，最后除以项数。它回答像素颜色差距，不直接回答三维结构或人的视觉偏好。
\par


\noindent\textbf{【人工教学：只有一个RGB像素。】}参考为(100,120,130)，预测为(110,100,130)，8位数值差为(10,-20,0)，平方和SSE=100+400+0=500。先按0至255的峰值归一化，MSE=500/(3×255²)=500/195075，约0.00256312。分母是3个颜色通道，不能误除以1个像素。
\par


\teaching{\textbf{正式主指标的分母固定，下面不是实测分数：}
    \[L_a=\frac14\sum_{t=20}^{23}\frac{1}{995328}
      \sum_{p,c}\left(\frac{U_{a,t,p,c}-U_{{\rm ref},t,p,c}}{255}\right)^2.\]
    差值 $\Delta=L_{\rm guide}-L_{\rm terminal}$；负号表示该MSE更低，不表示统计显著。}

假设另外三个教学画面的MSE分别是0.01、0.02、0.03，本页第一个画面仍为500/195075，四帧等权平均就是(500/195075+0.01+0.02+0.03)/4，约0.01564078。四个数均是教学构造，没有填入任何S86结果表。
\par


真实S86每帧576×576×3=995328个标量，全部四帧3981312个标量，每帧权重1/4。四臂均保留目标20、21、22、23，不能删掉某个很差的目标后继续称“四目标均值”。主指标使用原发图规则得到的uint8 RGB，除255后计算；独立复算可以int64累计SSE，避免把微小求和舍入当方向改变。
\par


主差值定义为Δ=L\_Gguide-L\_Gterminal，负值表示Gguide该指标较低；同时始终报告它相对G0和Gpaste的差。support与hole只能作为次级完整区域统计，区域空集写NA，不能写成0误差。旧目标参考已见过，预测封存后再读答案仍不恢复盲测。
\par


MSE下降可能奖励模糊。不能把分数低直接写成“看起来更真实”“几何更准确”，更不能用数百万通道标量冒充数百万独立场景。
\par


\clearpage

\section*{9. 结果出来后，哪些解释可以留下？}

情形一：Gguide不如G0。不能报总体改善；无论它是否比另一种融合好，都要保留原基线更好的事实。
\par


情形二：Gguide胜Gpaste，却不胜Gterminal。未观察到超越末步融合对照的额外MSE收益，不能据此确认某个机制已解释全部现象，也没有依据坚持较早融合更有益。
\par


情形三：Gguide胜Gterminal，但仍输G0。只能说两种干预中一个损害较小；不能把相对改善改写成超越原生成。
\par


情形四：Gguide同时胜过G0、Gpaste、Gterminal。可以保留“这个固定已见设置里，额外前49步干预及后续传播带来了像素误差收益”。还不能排除累计强度、简单边界处理或其他普通规则；仍需要独立场景和消融，才谈方法贡献。
\par


另有不完整情形：真实链超预算、缺目标、raw中有NaN/Inf、发图字节不一致或RNG核对失败。应保留真实失败与已产生结果，整体比较标为不完整或不可解释，不能补零、减分母或自动重跑直到成功。
\par


按本次合同，将保存Gguide全部50步raw/used clean，约66.36MB原始载荷；后验直接从保存量复算融合、末步Euler和发图量化，不再跑一条完整模型链来“复核”。旧S70同设置完整链约24.6分钟，S86两条链约50分钟是计划量级，不是当前实际完成时间；冻结总上限90分钟、每完整臂40分钟，资源限制失败仍属需报告的真实结果。
\par


本页没有任何真实S86分数。实际运行状态以文首截点及对应正式回执为准；本地文件出现或进度到第50步都不能替代完整终态、不同作者保存量复核和固定评分。
\par


\clearpage

\section*{10. 给导师的自然口述与自测}

可以这样讲：“我最近把历史几何真正推进成可供生成器使用的参考图。先用已有模型预测四张历史照片的几何，再固定相机做拟合，随后把历史颜色投到四个请求视角。投影阶段已完成，全部16对都保留，不同作者从原输入复算一致。预测覆盖从约94\%到79\%，但它只表示有历史候选，不是准确率。现在的关键不是给这套普通流程起新名字，而是看生成过程真的使用它以后有没有收益。因此我设置原生成、末端RGB合成、末步latent融合、多步融合四个对照。末步对照很重要，因为多步方法看起来更好，也可能只是最后融合与解码有效。真实生成状态和评分我会按完成回执报告；没有结果时不提前讲成功。”
\par


\noindent\textbf{自测1：为什么Z=2米不一定距离也是2米？}\ 答：Z是沿光轴分量；离光轴越远，欧氏距离sqrt(X²+Y²+Z²)越大。
\par


\noindent\textbf{自测2：一个源点有四个正足迹，是否新增四个独立观察？}\ 答：没有，是同一预测点对相邻格点的重复候选。
\par


\noindent\textbf{自测3：孔洞RGB为0，能否把所有黑像素都当洞？}\ 答：不能，真实黑色也可能为0，必须看独立mask。
\par


\noindent\textbf{自测4：为什么Gterminal不必重跑50步？}\ 答：它只消费同一G0真实末步保存量，做一次融合与原Euler末步算术，再解码。
\par


\noindent\textbf{自测5：Gguide只胜Gpaste可否证明较早引导有效？}\ 答：不能；latent域与解码本身是竞争解释，必须面对Gterminal。
\par


\noindent\textbf{自测6：复算最大差0能否证明世界几何正确？}\ 答：不能，两套程序可以把同一错误输入一致地算完。
\par


证据入口：S85\_RESULTS.md、S85的SUMMARY.json和\path{INDEPENDENT_OUTPUT_REVIEW.json}；S86的CONTRACT.json、ROOT\_LAUNCH\_01.json、\path{supervision_01/SUPERVISION.json}，以及完成后才出现的终态/复核/评分。旧142页报告仍保留其原科学截点，本增补不覆盖旧PDF。
\par


\vfill\noindent\small\textcolor{accent}{内容依据本地冻结协议与实际回执；人工教学单独标注。作者排版核验与root独立科学/视觉验收另见报告目录回执。}

\end{document}
